	\newpage
\section{Wnioski}	%5
%Npisać wnioski końcowe z przeprowadzonego projektu, 


NIS (Network Information Service) i NFS (Network File System) to protokoły sieciowe, które umożliwiają centralne zarządzanie informacjami o użytkownikach i grupach oraz udostępnianie systemów plików w sieci. W projekcie udało się zrealizować wszystkie założenia, a także zainstalować i skonfigurować usługi NIS i NFS na serwerze wirtualnym. Dzięki temu możliwe było automatyczne udostępnianie informacji o użytkownikach oraz współdzielenie katalogów domowych między różnymi systemami operacyjnymi w sieci.

Wszystkie usługi działają poprawnie, a użytkownicy mogą logować się do systemu i korzystać z udostępnionych zasobów. Projekt wykazał, że NIS i NFS są skutecznymi narzędziami do zarządzania informacjami o użytkownikach i współdzielenia plików w sieci. W przyszłości można rozważyć zastosowanie dodatkowych zabezpieczeń, takich jak szyfrowanie danych przesyłanych przez NIS, aby zwiększyć bezpieczeństwo systemu.
